\documentclass[11pt]{article}
\usepackage[margin=1in]{geometry}
\usepackage{amsmath,amssymb}
\usepackage{hyperref}

\title{Literature Status: The Real-Line Gaussian Tail-Space Subcase of Question 7}
\author{Run \texttt{fa\_banach\_001}, lane 1}
\date{2026-07-18}

\begin{document}
\maketitle

\noindent\textbf{Status.} Literature implied answer, real-line subcase. This is not a new proof packet.

\section*{Source Question}

The source paper is A. Eskenazis and P. Ivanisvili, \emph{Dimension independent Bernstein--Markov inequalities in Gauss space}, arXiv:1808.01273v3, Journal of Approximation Theory 253 (2020), Article 105377.

Question 7 asks whether, for every \(1<p<\infty\), there is \(c_p>0\) such that for every \(k\ge1\), every \(d\ge1\), and every polynomial
\[
P(x)=\sum_{|\alpha|\ge d} c_\alpha H_\alpha(x)
\]
on \(\mathbb R^k\), one has
\[
\|LP\|_{L^p(d\gamma_k)}\ge c_p d\,\|P\|_{L^p(d\gamma_k)}.
\]
Here \(L=\Delta-x\cdot\nabla\) is the Ornstein--Uhlenbeck generator, with \(LH_\alpha=-|\alpha|H_\alpha\) in the source convention.

\section*{Supporting Paper}

The supporting paper is A. Eskenazis and P. Ivanisvili, \emph{Sharp growth of the Ornstein--Uhlenbeck operator on Gaussian tail spaces}, arXiv:2011.01359v2; Israel Journal of Mathematics 253 (2023), 469--485, DOI \href{https://doi.org/10.1007/s11856-022-2373-8}{10.1007/s11856-022-2373-8}.

In the introduction, Corollary 2 states that for every \(p\in(1,\infty)\) there are constants \(S_p,T_p<\infty\) such that, for every one-dimensional polynomial \(f\) in the \(d\)-tail space,
\[
d\|f\|_{L^p(d\gamma_1)}\le T_p\|Lf\|_{L^p(d\gamma_1)}.
\]
Equivalently,
\[
\|Lf\|_{L^p(d\gamma_1)}\ge T_p^{-1}d\,\|f\|_{L^p(d\gamma_1)}.
\]
This is precisely the \(k=1\) case of Question 7.

\section*{Scope Limitation}

The supporting result answers only the real-line subcase of the source question. It does not settle Question 7 for arbitrary \(k\)-dimensional Hermite-tail polynomials, nor Question 3 on the sharp \(\sqrt{\deg P}\) gradient Bernstein--Markov inequality in all dimensions.

\section*{Provenance}

The supporting authors explicitly know the relation: arXiv:2011.01359 says the second inequality answers the Gaussian analogue of Mendel--Naor's question on the real line and cites the source paper arXiv:1808.01273. Because the original Question 7 is stated for all \(k\ge1\), this packet is classified as \(\texttt{literature\_implied\_answer (real-line subcase)}\), not as a full literature answer.

\end{document}
